% from ega4-18-fr.tex
\oldpage[IV-4]{120}

isomorphe à $A[T]/F.A[T]$, et pour tout autre polynôme unitaire $G\in A[T]$ tel que $B$ soit isomorphe à $A[T]/G.A[T]$, $G$ est nécessairement de degré $n$.

\begin{theorem}[18.4.6]
\label{IV.18.4.6}
(Chevalley). ---
\begin{enumerate}
  \item[(i)] Soient $f:X\to Y$ un morphisme localement de type fini, $x$ un point de $X$, $y=f(x)$, et posons $A=\sh{O}_{Y,y}$. Pour que $\sh{O}_{X,x}$ soit une $A$-algèbre formellement étale (resp. formellement non ramifiée), il faut et il suffit qu'il existe un polynôme unitaire $F\in A[T]$ et un idéal maximal $\mathfrak n$ de $B=A[T]/F.A[T]$ (resp. d'une algèbre quotient $B$ de $A[T]/F.A[T]$) tels que $\sh{O}_{X,x}$ soit $A$-isomorphe à $B_{\mathfrak n}$ et que, si $u$ est l'image de $T$ dans $B$, on ait $F'(u)\notin\mathfrak n$.
  \item[(ii)] Supposons en outre que $f$ soit localement de présentation finie. Alors, pour que $f$ soit étale au point $x$, il faut et il suffit de plus que l'on puisse prendre $B=A[T]/F.A[T]$.
\end{enumerate}
\end{theorem}

Les conditions sont suffisantes en vertu de (18.4.2). Pour voir qu'elles sont nécessaires, on peut évidemment se borner au cas où $X$ et $Y$ sont affines, et, compte tenu de la remarque (17.4.1.2), au cas où $f$ est formellement non ramifié et quasi-fini. Comme $f$ est affine, il résulte de (8.12.8) qu'il existe une $A$-algèbre finie $C$ et un idéal maximal $\mathfrak r$ de $C$ (nécessairement au-dessus de l'idéal maximal $\mathfrak m$ de $A$) tels que $\sh{O}_{X,x}$ soit $A$-isomorphe à $C_{\mathfrak r}$. En outre (17.4.1.2) le corps résiduel $C/\mathfrak r=C_{\mathfrak r}/\mathfrak rC_{\mathfrak r}$ est une extension finie séparable de $k=A/\mathfrak m$, donc de la forme $k[v]$, où $v$ est \emph{séparable} sur $k$. Soient $\mathfrak r_i$ ($1\leq i\leq h$) les idéaux maximaux de l'anneau semi-local $C$ \emph{autres} que $\mathfrak r$; il existe un élément $u\in C$ appartenant à tous les $\mathfrak r_i$ et tel que son image dans $C/\mathfrak r$ soit égale à $v$ (Bourbaki, \emph{Alg. comm.}, chap. II, \textsection~1, n\textsuperscript{o}~2, prop. 5). Nous allons montrer que la sous-$A$-algèbre $B=A[u]$ de $C$ et l'idéal (nécessairement maximal puisque $C$ est finie sur $B$) $\mathfrak n=\mathfrak r\cap B$ de $B$ répondent à la question.

Pour traiter le cas où $\sh{O}_{X,x}$ est une $A$-algèbre formellement ramifiée, il suffira de prouver que $B_{\mathfrak n}$ est isomorphe à $C_{\mathfrak r}$; en effet, $B_{\mathfrak n}$ sera alors formellement non ramifié sur $A$ et l'existence du polynôme $F\in A[T]$ ayant les propriétés de l'énoncé résultera de (18.4.4). Notons maintenant que, puisque $f$ est formellement non ramifié, $C/\mathfrak r$ est isomorphe à la $A$-algèbre $C_{\mathfrak r}/\mathfrak mC_{\mathfrak r}$ (17.4.1.2). On est ainsi ramené à prouver le lemme suivant:

\begin{lemma}[18.4.6.1]
\label{IV.18.4.6.1}
Soient $A$ un anneau local, $\mathfrak m$ son idéal maximal, $C$ une $A$-algèbre finie, $\mathfrak r$ un idéal maximal de $C$. Soit $u$ un élément de $C$ appartenant à tous les idéaux maximaux de $C$ distincts de $\mathfrak r$, n'appartenant pas à $\mathfrak r$, et tel que $C_{\mathfrak r}/\mathfrak mC_{\mathfrak r}$ soit une algèbre monogène sur $k=A/\mathfrak m$, engendrée par l'image de $u$ dans $C_{\mathfrak r}/\mathfrak mC_{\mathfrak r}$. Posons $B=A[u]$, $\mathfrak n=\mathfrak r\cap B$. Alors l'homomorphisme canonique $B_{\mathfrak n}\to C_{\mathfrak r}$ est un isomorphisme.
\end{lemma}

Posons $R=B-\mathfrak n$, $S=C-\mathfrak r$, de sorte que $B_{\mathfrak n}=R^{-1}B$ et $C_{\mathfrak r}=S^{-1}C$; l'homomorphisme canonique $B_{\mathfrak n}\to C_{\mathfrak r}$ peut s'écrire comme le composé
\[
  R^{-1}B\xrightarrow{g}R^{-1}C\xrightarrow{h}S^{-1}C
\]
et il suffit de montrer que chacun de ces deux homomorphismes est bijectif.

Montrons d'abord que $h:R^{-1}C\to S^{-1}C=C_{\mathfrak r}$ est bijectif; il suffit de voir que les images dans $R^{-1}C$ des éléments de $S$ sont inversibles, ou encore que tout idéal maximal $\mathfrak p$ de $R^{-1}C$ a une image réciproque dans $C$ ne rencontrant pas $S$, donc nécessairement égale à $\mathfrak r$. Or, comme $R^{-1}C$ est une $(R^{-1}B)$-algèbre finie, l'image réciproque de $\mathfrak p$ dans

